(1) 固定电荷时真空电场能 $U_Q=Q^2/[2C(x)]=Q^2x/(2\varepsilon_0A)$。取 $x$ 正向表示板间距增加，则保守力 $F_x=-\partial U_Q/\partial x=-Q^2/(2\varepsilon_0A)<0$，确为相互吸引，其大小在理想无限平板固定电荷模型中与间隙无关。
(2) 恒压时电源不断供应电荷，电场能 $U_{\rm field}=C(x)V^2/2$，电源自由能变化对应 Legendre 变换 $\mathcal F_V=U_{\rm field}-VQ=-C(x)V^2/2=-\varepsilon_0AV^2/(2x)$。于是 $F_x=-\partial\mathcal F_V/\partial x=-\varepsilon_0AV^2/(2x^2)$，若对 $U_{\rm field}$ 直接求导，会错误预测排斥。
(3) 总恒压有效势 $\mathcal F(x)=\frac12k(x-x_0)^2-\varepsilon_0AV^2/(2x)$。平衡条件 $k(x-x_0)+\varepsilon_0AV^2/(2x^2)=0$，稳定条件为二阶导 $\mathcal F''(x)=k-\varepsilon_0AV^2/x^3>0$；如果二阶导负则是失稳临界后的非稳定平衡。
(4) 拉入点同时满足 $k(x_*-x_0)+\varepsilon_0AV_*^2/(2x_*^2)=0$ 与 $k=\varepsilon_0AV_*^2/x_*^3$。消去 $V_*$ 得 $x_*=2x_0/3$、$V_*^2=8kx_0^3/(27\varepsilon_0A)$；两式具有长度和电压平方的正确单位。近接触时真实粗糙度、弹簧非线性与边缘场将改变精确阈值。
